\subsection{约束问题}
我们常常遇到的一类问题是$u$满足其他的约束条件,比如上述两点之间曲线长度问题可以增加一个约束条件,
所构成的曲线所围面积为定值$A$,即$\int u(x) dx = A$.通常的做法是利用拉格朗日乘子法,即对
目标函数$J[u]$添加一项$\lambda\left( \int u(x) dx - A \right)$,变成无约束问题去求解. 
这里的$\lambda$称为\textbf{拉格朗日乘子}.
以上例为例,我们需要极小化的函数变成
$$
J[u] = \int_{x_1}^{x_2} \sqrt{ 1 + u'(x)^2} dx   + \lambda \left( \int_{x_1}^{x_2} u(x) dx - A \right).
$$
这样拉式量内包含了这一约束.对于该例子应用欧拉-拉格朗日方程,我们有
\[
  \lambda - \funcd{}{x} \frac{u'}{\sqrt{1 + u'^2}} = 0,
\]
可以求得
\[
 \frac{u'}{\sqrt{1+u'^2}}= \lambda x + c   %\textrm{或} u'=\frac{\lambda x-c}{\sqrt{1-(\lambda x-c)^2}}
 \]
令辅助函数$ p(x)\quad  = \lambda x + c   $,  两边平方,得到
\[ u' = \frac{p}{\sqrt{1 - p^2}} .
\]
换元令$ t = \lambda x + c$, 则有 $dt = \lambda dx$, 于是
\[
 du = u' dx = \frac{t}{\sqrt{1 - t^2}} \frac{dt}{\lambda} .
\]
积分得到
\[
u(x) = \int \frac{t}{\sqrt{1 - t^2}} \frac{dt}{\lambda} = - \frac{\sqrt{1 - t^2}}{\lambda} + d,
\]
其中$d$为积分常数.
最终得到方程
\[
  (\lambda x +c)^2 + (\lambda u -d)^2 = 1
\]
其中$c,d,\lambda$由$u(x_1)=y_1, u(x_2) = y_2$ 和$\int u(x) dx = A$决定.
在这样的约束条件下,最短路径是一段圆弧！
此外, 可以对$p(x)$直接求导得到
\[
\lambda = \frac{u''}{(1+y'^2) ^{3/2}}
\]
而右边正式平面曲线的曲率$\kappa$, 因此在约束条件下, 最短路径的曲率是常数.
故也可知曲线为圆弧.

\begin{example}
悬链线.
\end{example}
\begin{solution}
均匀软绳的重力势能是
\begin{equation}
  U[u] = \int_{-L/2}^{L/2} \rho g\, u\sqrt{1+u'^2}\,\mathrm{d}x,
\end{equation}
绳长另有约束 $\int_{-L/2}^{L/2}\sqrt{1+u'^2}\,\mathrm{d}x=\ell$.
乘子 $\lambda$ 把被积函数改成 $(\rho g u+\lambda)\sqrt{1+u'^2}$.
它不显含 $x$,首次积分给出
\begin{equation}
  \frac{\rho g u+\lambda}{\sqrt{1+u'^2}} = A .
\end{equation}
令 $w=\rho g u+\lambda$,则 $\mathrm{d}w/\mathrm{d}x = (\rho g/A)\sqrt{w^2-A^2}$.
分离变量并积分,
\begin{equation}
  w = A \cosh\Bigl(\frac{\rho g}{A}(x-x_c)\Bigr).
\end{equation}
因此
\begin{equation}
  u(x) = a\cosh\Bigl(\frac{x-x_c}{a}\Bigr) + u_c,
  \qquad a=\frac{A}{\rho g}.
\end{equation}
平衡曲线是双曲余弦,不是抛物线.顶点附近把余弦展开到二阶,才近似成抛物线.常数由两端位置和绳长决定.
\end{solution}


% \subsection{Euler-Lagrange方程}
% 上式的求解其实可以归结

% \subsubsection{二维问题}
